\subsection{$q$-binomial coefficients}

Next, we shall study $q$\emph{-binomial coefficients} (also known as
\emph{Gaussian binomial coefficients}, due to their origins in Gauss's
number-theoretical research \cite[\S 5]{Gauss08}). While we will define them
as generating functions for certain kinds of partitions, they are sufficiently
elementary to have relevance to various other subjects. We will scratch the
surface; more can be found in \cite[Chapters 5--7]{KacChe02}, \cite[Chapter
2]{KliSch97}, \cite[Chapter 5]{Wagner08}, \cite[\S 11.3--11.11]{Wagner20},
\cite[spread across the text]{Stanley-EC1}, \cite[\S 2.6]{GouJac83},
\cite{Johnso20} and other sources. (The book \cite{Johnso20} is particularly
recommended as a leisurely introduction to $q$-binomial coefficients and
related power series.)

\subsubsection{\label{subsec.qbin.motiv}Motivation}

For any $n\in\mathbb{N}$, we have%
\[
p\left(  n\right)  =\left(  \text{\# of partitions of }n\right)  .
\]
For any $n,k\in\mathbb{N}$, we have%
\begin{align*}
p_{k}\left(  n\right)   &  =\left(  \text{\# of partitions of }n\text{ into
}k\text{ parts}\right) \\
&  =\left(  \text{\# of partitions of }n\text{ with largest part }k\right)
\ \ \ \ \ \ \ \ \ \ \left(  \text{by Theorem \ref{prop.pars.pkn=dual}}\right)
.
\end{align*}
Thus, for any $n,k\in\mathbb{N}$, we have%
\begin{align*}
p_{0}\left(  n\right)  +p_{1}\left(  n\right)  +\cdots+p_{k}\left(  n\right)
&  =\left(  \text{\# of partitions of }n\text{ into \textbf{at most} }k\text{
parts}\right) \\
&  =\left(  \text{\# of partitions of }n\text{ with largest part }\leq
k\right)
\end{align*}
(by Corollary \ref{cor.pars.p0kn=dual}).

So far, so good. But how to count partitions of $n$ that both have a fixed \#
of parts (say, $k$ parts) and a fixed largest part (say, largest part $\ell$) ?

Let us first drop the size requirement -- i.e., we replace \textquotedblleft
partitions of $n$\textquotedblright\ by just \textquotedblleft
partitions\textquotedblright.

For example, how many partitions have $4$ parts and largest part $6$ ?

As in the proof of Theorem \ref{prop.pars.pkn=dual}, let us draw the Young
diagram of such a partition: For example, the partition $\left(
6,3,3,2\right)  $ has Young diagram%
\[
\ydiagram{6,3,3,2}\ \ .
\]


Consider the lower boundary of this Young diagram -- i.e., the
\textquotedblleft irregular\textquotedblright\ southeastern border between
what is in the diagram and what is outside of it. Let me mark it in thick red:%
\[%
% [tikz figure removed]
%BeginExpansion
% [tikz figure removed]%
%EndExpansion
\ \ .
\]
This lower boundary can be viewed as a lattice path from the point $\left(
0,0\right)  $ to the point $\left(  6,4\right)  $ (where we are using
Cartesian coordinates to label the intersections of grid lines, so that the
southwesternmost point in our diagram is $\left(  0,0\right)  $; note that
this is completely unrelated to our labeling of cells used in defining the
Young diagram!\footnote{For additional clarity, here are the Cartesian
coordinates of all grid points on our lattice path:%
\[%
%TCIMACRO{\TeXButton{tikz diagram with border and coordinates}%
%{% [tikz figure removed]}}%
%BeginExpansion
% [tikz figure removed]%
%EndExpansion
\ \ .
\]
}). This lattice path consists of east-steps (i.e., steps $\left(  i,j\right)
\rightarrow\left(  i+1,j\right)  $) and north-steps (i.e., steps $\left(
i,j\right)  \rightarrow\left(  i,j+1\right)  $); moreover, it begins with an
east-step (since otherwise, our partition would have fewer than $4$ parts) and
ends with a north-step (since otherwise, our partition would have largest part
$<6$). Moreover, the Young diagram (and thus the partition) is uniquely
determined by this lattice path, since its cells are precisely the cells
\textquotedblleft northwest\textquotedblright\ of this lattice path.
Conversely, any lattice path from $\left(  0,0\right)  $ to $\left(
6,4\right)  $ that consists of east-steps and north-steps and begins with an
east step and ends with a north-step uniquely determines a Young diagram and
therefore a partition. Therefore, in order to count the partitions that have
$4$ parts and largest part $6$, we only need to count such lattice paths.

To count them, we notice that any such lattice path has precisely $10$ steps
(since any step increases the sum of the coordinates by $1$; but this sum must
increase from $0+0=0$ to $6+4=10$). The first and the last steps are
predetermined; it remains to decide which of the remaining $8$ steps are
north-steps. The \# of ways to do this is $\dbinom{8}{3}$, because we want
precisely $3$ of our $8$ non-predetermined steps to be north-steps (in order
to end up at $\left(  6,4\right)  $ rather than some other point).

As a consequence of this all, we find%
\[
\left(  \text{\# of partitions with }4\text{ parts and largest part }6\right)
=\dbinom{8}{3}.
\]
More generally, by the same argument, we obtain the following:

\begin{proposition}
\label{prop.pars.qbinom.intro-count-binom}For any positive integers $k$ and
$\ell$, we have%
\[
\left(  \text{\# of partitions with }k\text{ parts and largest part }%
\ell\right)  =\dbinom{k+\ell-2}{k-1}.
\]

\end{proposition}

Note two things:

\begin{itemize}
\item This is a finite number, even without fixing the size of the partition.
This is not surprising, since you have only finitely many parts and only
finitely many options for each part.

\item The number is symmetric in $k$ and $\ell$. This, too, is not surprising,
because conjugation (as defined in the proof of Theorem
\ref{prop.pars.pkn=dual}) gives a bijection%
\begin{align*}
&  \text{from }\left\{  \text{partitions with }k\text{ parts and largest part
}\ell\right\} \\
&  \text{to }\left\{  \text{partitions with }\ell\text{ parts and largest part
}k\right\}  .
\end{align*}

\end{itemize}

Now, let us integrate the size of the partition back into our count -- i.e.,
let us try to count the partitions of a given $n\in\mathbb{N}$ with $k$ parts
and largest part $\ell$. No simple formula (like Proposition
\ref{prop.pars.qbinom.intro-count-binom}) exists for this number any more, so
we switch our focus to the generating function of such numbers (for fixed $k$
and $\ell$). In other words, we try to compute the FPS%
\begin{align*}
&  \sum_{n\in\mathbb{N}}\left(  \text{\# of partitions of }n\text{ with
}k\text{ parts and largest part }\ell\right)  x^{n}\\
&  =\sum_{\substack{\lambda\text{ is a partition}\\\text{with largest part
}\ell\\\text{and length }k}}x^{\left\vert \lambda\right\vert }.
\end{align*}


For reasons of convenience, history and simplicity, we modify this problem
slightly (without changing its essence). To wit,

\begin{itemize}
\item we rename $\ell$ as $n-k$ (note that $n$ will no longer stand for the
size of the partition);

\item we replace \textquotedblleft largest part $n-k$ and length
$k$\textquotedblright\ by \textquotedblleft largest part $\leq n-k$ and length
$\leq k$\textquotedblright\ (this changes the results of our counts, but we
can easily recover the answer to the original question from an answer to the
new question; e.g., in order to count the length-$k$ partitions, it suffices
to subtract the \# of length-$\left(  \leq k-1\right)  $-partitions from the
\# of length-$\left(  \leq k\right)  $ partitions);

\item we rename the indeterminate $x$ as $q$.
\end{itemize}

\subsubsection{Definition}

\begin{convention}
\label{conv.pars.qbinom.q}In this section, we will mostly be using FPSs in the
indeterminate $q$. That is, we call the indeterminate $q$ rather than $x$.
Thus, e.g., our formula%
\[
\prod_{n>0}\left(  1-x^{n}\right)  ^{-1}=\prod_{n>0}\dfrac{1}{1-x^{n}}%
=\sum_{n\in\mathbb{N}}p\left(  n\right)  x^{n}=\sum_{\substack{\lambda\text{
is a}\\\text{partition}}}x^{\left\vert \lambda\right\vert }%
\]
becomes%
\[
\prod_{n>0}\left(  1-q^{n}\right)  ^{-1}=\prod_{n>0}\dfrac{1}{1-q^{n}}%
=\sum_{n\in\mathbb{N}}p\left(  n\right)  q^{n}=\sum_{\substack{\lambda\text{
is a}\\\text{partition}}}q^{\left\vert \lambda\right\vert }.
\]
The ring of FPSs in the indeterminate $q$ over a commutative ring $K$ will be
denoted by $K\left[  \left[  q\right]  \right]  $. The ring of polynomials in
the indeterminate $q$ over $K$ will be denoted by $K\left[  q\right]  $.
\end{convention}

\begin{definition}
\label{def.pars.qbinom.qbinom}Let $n\in\mathbb{N}$ and $k\in\mathbb{Z}$.
\medskip

\textbf{(a)} The $q$\emph{-binomial coefficient} (or \emph{Gaussian binomial
coefficient}) $\dbinom{n}{k}_{q}$ is defined to be the polynomial%
\[
\sum_{\substack{\lambda\text{ is a partition}\\\text{with largest part }\leq
n-k\\\text{and length }\leq k}}q^{\left\vert \lambda\right\vert }\in
\mathbb{Z}\left[  q\right]  .
\]


This is also denoted by $%
%TCIMACRO{\QDATOPD{[}{]}{n}{k}}%
%BeginExpansion
\genfrac{[}{]}{0pt}{0}{n}{k}%
%EndExpansion
$ (but this notation has other meanings, too, and suppresses $q$). \medskip

\textbf{(b)} If $a$ is any element of a ring $A$, then we set%
\begin{align*}
\dbinom{n}{k}_{a}:=  &  \ \dbinom{n}{k}_{q}\left[  a\right] \\
&  \ \ \ \ \ \ \ \ \ \ \left(  \text{this means the result of substituting
}a\text{ for }q\text{ in }\dbinom{n}{k}_{q}\right) \\
=  &  \ \sum_{\substack{\lambda\text{ is a partition}\\\text{with largest part
}\leq n-k\\\text{and length }\leq k}}a^{\left\vert \lambda\right\vert }\in A.
\end{align*}

\end{definition}

\begin{remark}
\label{rmk.pars.qbinom.poly}The $\dbinom{n}{k}_{q}$ we defined in Definition
\ref{def.pars.qbinom.qbinom} \textbf{(a)} is really a polynomial, not merely a
FPS, because (for any given $n$ and $k$) there are only finitely many
partitions with largest part $\leq n-k$ and length $\leq k$.
\end{remark}

\begin{remark}
\label{rmk.pars.qbinom.neg}The notation $\dbinom{n}{k}_{q}$ (and the name
\textquotedblleft$q$-binomial coefficient\textquotedblright) suggests a
similarity to the usual binomial coefficient $\dbinom{n}{k}$. And indeed, we
will soon see that $\dbinom{n}{k}_{1}=\dbinom{n}{k}$.

Note, however, that $\dbinom{n}{k}_{q}$ is only defined for $n\in\mathbb{N}$
and $k\in\mathbb{Z}$ (unlike $\dbinom{n}{k}$, which we defined for arbitrary
$n,k\in\mathbb{C}$). It is possible to extend it to negative integers $n$, but
this will result in a Laurent polynomial. (See Exercise
\ref{exe.pars.qbinom.upneg} for this extension.)
\end{remark}

\begin{example}
We have%
\begin{align*}
\dbinom{3}{2}_{q}  &  =\sum_{\substack{\lambda\text{ is a partition}%
\\\text{with largest part }\leq1\\\text{and length }\leq2}}q^{\left\vert
\lambda\right\vert }=q^{\left\vert \left(  1,1\right)  \right\vert
}+q^{\left\vert \left(  1\right)  \right\vert }+q^{\left\vert \left(
{}\right)  \right\vert }\\
&  \ \ \ \ \ \ \ \ \ \ \ \ \ \ \ \ \ \ \ \ \left(
\begin{array}
[c]{c}%
\text{since the partitions with largest part }\leq1\\
\text{and length }\leq2\text{ are }\left(  1,1\right)  \text{, }\left(
1\right)  \text{ and }\left(  {}\right)
\end{array}
\right) \\
&  =q^{2}+q^{1}+q^{0}=q^{2}+q+1
\end{align*}
and%
\begin{align*}
\dbinom{4}{2}_{q}  &  =\sum_{\substack{\lambda\text{ is a partition}%
\\\text{with largest part }\leq2\\\text{and length }\leq2}}q^{\left\vert
\lambda\right\vert }\\
&  =q^{\left\vert \left(  2,2\right)  \right\vert }+q^{\left\vert \left(
2,1\right)  \right\vert }+q^{\left\vert \left(  2\right)  \right\vert
}+q^{\left\vert \left(  1,1\right)  \right\vert }+q^{\left\vert \left(
1\right)  \right\vert }+q^{\left\vert \left(  {}\right)  \right\vert }\\
&  =q^{4}+q^{3}+q^{2}+q^{2}+q^{1}+q^{0}=q^{4}+q^{3}+2q^{2}+q+1.
\end{align*}

\end{example}

\subsubsection{Basic properties}

Let us show two slightly different (but equivalent) ways to express
$q$-binomial coefficients:

\begin{proposition}
\label{prop.pars.qbinom.alt-defs}Let $n\in\mathbb{N}$ and $k\in\mathbb{Z}$.
\medskip

\textbf{(a)} We have%
\[
\dbinom{n}{k}_{q}=\sum_{0\leq i_{1}\leq i_{2}\leq\cdots\leq i_{k}\leq
n-k}q^{i_{1}+i_{2}+\cdots+i_{k}}.
\]
Here, the sum ranges over all weakly increasing $k$-tuples $\left(
i_{1},i_{2},\ldots,i_{k}\right)  \in\left\{  0,1,\ldots,n-k\right\}  ^{k}$. If
$k>n$, then this is an empty sum (since the set $\left\{  0,1,\ldots
,n-k\right\}  $ is empty in this case, and thus its $k$-th power $\left\{
0,1,\ldots,n-k\right\}  ^{k}$ is also empty because $k>n\geq0$). If $k<0$,
then this is also an empty sum (since there are no $k$-tuples when $k<0$).
\medskip

\textbf{(b)} Set $\operatorname*{sum}S=\sum_{s\in S}s$ for any finite set $S$
of integers. (For example, $\operatorname*{sum}\left\{  2,4,5\right\}
=2+4+5=11$.) Then, we have%
\[
\dbinom{n}{k}_{q}=\sum_{\substack{S\subseteq\left\{  1,2,\ldots,n\right\}
;\\\left\vert S\right\vert =k}}q^{\operatorname*{sum}S-\left(  1+2+\cdots
+k\right)  }.
\]


\textbf{(c)} We have%
\[
\dbinom{n}{k}_{1}=\dbinom{n}{k}.
\]

\end{proposition}

\begin{example}
For example, let us compute $\dbinom{5}{2}_{q}$ using Proposition
\ref{prop.pars.qbinom.alt-defs} \textbf{(b)}. Namely, applying Proposition
\ref{prop.pars.qbinom.alt-defs} \textbf{(b)} to $n=5$ and $k=2$, we obtain%
\begin{align*}
\dbinom{5}{2}_{q}  &  =\sum_{\substack{S\subseteq\left\{  1,2,\ldots
,5\right\}  ;\\\left\vert S\right\vert =2}}q^{\operatorname*{sum}S-\left(
1+2\right)  }\\
&  =q^{\left(  1+2\right)  -\left(  1+2\right)  }+q^{\left(  1+3\right)
-\left(  1+2\right)  }+q^{\left(  1+4\right)  -\left(  1+2\right)
}+q^{\left(  1+5\right)  -\left(  1+2\right)  }\\
&  \ \ \ \ \ \ \ \ \ \ +q^{\left(  2+3\right)  -\left(  1+2\right)
}+q^{\left(  2+4\right)  -\left(  1+2\right)  }+q^{\left(  2+5\right)
-\left(  1+2\right)  }\\
&  \ \ \ \ \ \ \ \ \ \ +q^{\left(  3+4\right)  -\left(  1+2\right)
}+q^{\left(  3+5\right)  -\left(  1+2\right)  }+q^{\left(  4+5\right)
-\left(  1+2\right)  }\\
&  \ \ \ \ \ \ \ \ \ \ \ \ \ \ \ \ \ \ \ \ \left(
\begin{array}
[c]{c}%
\text{since the }2\text{-element subsets of }\left\{  1,2,\ldots,5\right\}
\text{ are}\\
\left\{  1,2\right\}  ,\ \left\{  1,3\right\}  ,\ \left\{  1,4\right\}
,\ \left\{  1,5\right\}  ,\ \left\{  2,3\right\}  ,\ \left\{  2,4\right\}  ,\\
\left\{  2,5\right\}  ,\ \left\{  3,4\right\}  ,\ \left\{  3,5\right\}
,\ \left\{  4,5\right\}
\end{array}
\right) \\
&  =q^{0}+q^{1}+q^{2}+q^{3}+q^{2}+q^{3}+q^{4}+q^{4}+q^{5}+q^{6}\\
&  =1+q+2q^{2}+2q^{3}+2q^{4}+q^{5}+q^{6}.
\end{align*}

\end{example}

\begin{proof}
[Proof of Proposition \ref{prop.pars.qbinom.alt-defs}.]\textbf{(a)} The
definition of $\dbinom{n}{k}_{q}$ yields%
\begin{align}
\dbinom{n}{k}_{q}  &  =\sum_{\substack{\lambda\text{ is a partition}%
\\\text{with largest part }\leq n-k\\\text{and length }\leq k}}q^{\left\vert
\lambda\right\vert }\nonumber\\
&  =\sum_{\ell=0}^{k}\ \ \sum_{\substack{\lambda\text{ is a partition}%
\\\text{with largest part }\leq n-k\\\text{and length }\ell}}q^{\left\vert
\lambda\right\vert }. \label{pf.prop.pars.qbinom.alt-defs.a.1}%
\end{align}
Now, let us simplify the inner sum on the right hand side.

Fix $\ell\in\left\{  0,1,\ldots,k\right\}  $. Then, any partition $\lambda$
with length $\ell$ has the form $\left(  \lambda_{1},\lambda_{2}%
,\ldots,\lambda_{\ell}\right)  $ for some nonnegative integers $\lambda
_{1},\lambda_{2},\ldots,\lambda_{\ell}$ satisfying $\lambda_{1}\geq\lambda
_{2}\geq\cdots\geq\lambda_{\ell}>0$ (by the definitions of \textquotedblleft
partition\textquotedblright\ and \textquotedblleft length\textquotedblright).
Moreover, this partition $\lambda$ has largest part $\leq n-k$ if and only if
its entries satisfy $n-k\geq\lambda_{1}\geq\lambda_{2}\geq\cdots\geq
\lambda_{\ell}>0$. Finally, the size $\left\vert \lambda\right\vert $ of this
partition equals $\lambda_{1}+\lambda_{2}+\cdots+\lambda_{\ell}$. Hence, we
can rewrite the sum%
\[
\sum_{\substack{\lambda\text{ is a partition}\\\text{with largest part }\leq
n-k\\\text{and length }\ell}}q^{\left\vert \lambda\right\vert }%
\ \ \ \ \ \ \ \ \ \ \text{as}\ \ \ \ \ \ \ \ \ \ \sum_{\substack{\left(
\lambda_{1},\lambda_{2},\ldots,\lambda_{\ell}\right)  \in\mathbb{N}^{\ell
};\\n-k\geq\lambda_{1}\geq\lambda_{2}\geq\cdots\geq\lambda_{\ell}%
>0}}q^{\lambda_{1}+\lambda_{2}+\cdots+\lambda_{\ell}}.
\]
In other words, we have%
\begin{equation}
\sum_{\substack{\lambda\text{ is a partition}\\\text{with largest part }\leq
n-k\\\text{and length }\ell}}q^{\left\vert \lambda\right\vert }=\sum
_{\substack{\left(  \lambda_{1},\lambda_{2},\ldots,\lambda_{\ell}\right)
\in\mathbb{N}^{\ell};\\n-k\geq\lambda_{1}\geq\lambda_{2}\geq\cdots\geq
\lambda_{\ell}>0}}q^{\lambda_{1}+\lambda_{2}+\cdots+\lambda_{\ell}}.
\label{pf.prop.pars.qbinom.alt-defs.a.2}%
\end{equation}


Next, for any $k$-tuple $\left(  \lambda_{1},\lambda_{2},\ldots,\lambda
_{k}\right)  \in\mathbb{N}^{k}$, let us define $\operatorname*{numpos}\left(
\lambda_{1},\lambda_{2},\ldots,\lambda_{k}\right)  $ to be the number of
positive entries of this $k$-tuple (i.e., the number of $i\in\left\{
1,2,\ldots,k\right\}  $ satisfying $\lambda_{i}>0$). For example,
$\operatorname*{numpos}\left(  4,2,2,0\right)  =3$ and \newline%
$\operatorname*{numpos}\left(  4,2,2,1\right)  =4$ and $\operatorname*{numpos}%
\left(  0,0,0,0\right)  =0$. The following is easy but important:

\begin{statement}
\textit{Observation 1:} Let $\left(  \lambda_{1},\lambda_{2},\ldots
,\lambda_{k}\right)  \in\mathbb{N}^{k}$ be any $k$-tuple satisfying
$\lambda_{1}\geq\lambda_{2}\geq\cdots\geq\lambda_{k}\geq0$ and
$\operatorname*{numpos}\left(  \lambda_{1},\lambda_{2},\ldots,\lambda
_{k}\right)  =\ell$. Then:

\textbf{(a)} The first $\ell$ entries of $\left(  \lambda_{1},\lambda
_{2},\ldots,\lambda_{k}\right)  $ are positive (i.e., we have $\lambda_{i}>0$
for all $i\in\left\{  1,2,\ldots,\ell\right\}  $).

\textbf{(b)} The last $k-\ell$ entries of $\left(  \lambda_{1},\lambda
_{2},\ldots,\lambda_{k}\right)  $ are $0$ (i.e., we have $\lambda_{i}=0$ for
all $i\in\left\{  \ell+1,\ell+2,\ldots,k\right\}  $).

\textbf{(c)} We have $\lambda_{1}+\lambda_{2}+\cdots+\lambda_{\ell}%
=\lambda_{1}+\lambda_{2}+\cdots+\lambda_{k}$.

\textbf{(d)} The $\ell$-tuple $\left(  \lambda_{1},\lambda_{2},\ldots
,\lambda_{\ell}\right)  $ is a partition.
\end{statement}

\begin{fineprint}
[\textit{Proof of Observation 1:} We have $\operatorname*{numpos}\left(
\lambda_{1},\lambda_{2},\ldots,\lambda_{k}\right)  =\ell$. In other words, the
$k$-tuple $\left(  \lambda_{1},\lambda_{2},\ldots,\lambda_{k}\right)  $ has
exactly $\ell$ positive entries. Since this $k$-tuple is weakly decreasing
(because $\lambda_{1}\geq\lambda_{2}\geq\cdots\geq\lambda_{k}$), these $\ell$
positive entries must be concentrated at the beginning of the $k$-tuple; i.e.,
they must be the first $\ell$ entries of the $k$-tuple. Hence, the first
$\ell$ entries of $\left(  \lambda_{1},\lambda_{2},\ldots,\lambda_{k}\right)
$ are positive. This proves Observation 1 \textbf{(a)}.

We have shown that the $k$-tuple $\left(  \lambda_{1},\lambda_{2}%
,\ldots,\lambda_{k}\right)  $ has exactly $\ell$ positive entries, and they
are the first $\ell$ entries of this $k$-tuple. Hence, the remaining $k-\ell$
entries of this $k$-tuple are nonpositive. Since these entries are nonnegative
as well (because $\lambda_{1}\geq\lambda_{2}\geq\cdots\geq\lambda_{k}\geq0$),
we thus conclude that they are $0$. In other words, the last $k-\ell$ entries
of $\left(  \lambda_{1},\lambda_{2},\ldots,\lambda_{k}\right)  $ are $0$. This
proves Observation 1 \textbf{(b)}.

Furthermore,%
\begin{align*}
\lambda_{1}+\lambda_{2}+\cdots+\lambda_{k}  &  =\left(  \lambda_{1}%
+\lambda_{2}+\cdots+\lambda_{\ell}\right)  +\underbrace{\left(  \lambda
_{\ell+1}+\lambda_{\ell+2}+\cdots+\lambda_{k}\right)  }_{\substack{=0+0+\cdots
+0\\\text{(by Observation 1 \textbf{(b)})}}}\\
&  =\left(  \lambda_{1}+\lambda_{2}+\cdots+\lambda_{\ell}\right)  +\left(
0+0+\cdots+0\right)  =\lambda_{1}+\lambda_{2}+\cdots+\lambda_{\ell}.
\end{align*}
This proves Observation 1 \textbf{(c)}.

\textbf{(d)} The $\ell$-tuple $\left(  \lambda_{1},\lambda_{2},\ldots
,\lambda_{\ell}\right)  $ is weakly decreasing (since $\lambda_{1}\geq
\lambda_{2}\geq\cdots\geq\lambda_{k}$ entails $\lambda_{1}\geq\lambda_{2}%
\geq\cdots\geq\lambda_{\ell}$) and consists of positive integers (since
Observation 1 \textbf{(a)} says that we have $\lambda_{i}>0$ for all
$i\in\left\{  1,2,\ldots,\ell\right\}  $). Thus, it is a weakly decreasing
tuple of positive integers, i.e., a partition. This proves Observation 1
\textbf{(d)}.] \medskip
\end{fineprint}

Let us recall that $\ell\in\left\{  0,1,\ldots,k\right\}  $, so that $\ell\leq
k$. Hence, any $\ell$-tuple $\left(  \lambda_{1},\lambda_{2},\ldots
,\lambda_{\ell}\right)  \in\mathbb{N}^{\ell}$ can be extended to a $k$-tuple
$\left(  \lambda_{1},\lambda_{2},\ldots,\lambda_{k}\right)  \in\mathbb{N}^{k}$
by inserting $k-\ell$ zeroes at the end (i.e., by setting $\lambda_{\ell
+1}=\lambda_{\ell+2}=\cdots=\lambda_{k}=0$).\ \ \ \ \footnote{For example, if
$\ell=3$ and $k=5$, then the $\ell$-tuple $\left(  4,2,2\right)  $ gets
extended to the $k$-tuple $\left(  4,2,2,0,0\right)  $.} If the original
$\ell$-tuple $\left(  \lambda_{1},\lambda_{2},\ldots,\lambda_{\ell}\right)
\in\mathbb{N}^{\ell}$ was a partition with largest part $\leq n-k$, then the
extended $k$-tuple $\left(  \lambda_{1},\lambda_{2},\ldots,\lambda_{k}\right)
=\left(  \lambda_{1},\lambda_{2},\ldots,\lambda_{\ell},\underbrace{0,0,\ldots
,0}_{k-\ell\text{ zeroes}}\right)  $ will satisfy $n-k\geq\lambda_{1}%
\geq\lambda_{2}\geq\cdots\geq\lambda_{k}\geq0$ (since $n-k\geq\lambda_{1}%
\geq\lambda_{2}\geq\cdots\geq\lambda_{\ell}$ and $\lambda_{\ell}\geq
0=\lambda_{\ell+1}=\lambda_{\ell+2}=\cdots=\lambda_{k}\geq0$) and
$\operatorname*{numpos}\left(  \lambda_{1},\lambda_{2},\ldots,\lambda
_{k}\right)  =\ell$ (since its first $\ell$ entries are positive whereas its
remaining $k-\ell$ entries are $0$).

Conversely, if $\left(  \lambda_{1},\lambda_{2},\ldots,\lambda_{k}\right)
\in\mathbb{N}^{k}$ is a $k$-tuple satisfying $n-k\geq\lambda_{1}\geq
\lambda_{2}\geq\cdots\geq\lambda_{k}\geq0$ and $\operatorname*{numpos}\left(
\lambda_{1},\lambda_{2},\ldots,\lambda_{k}\right)  =\ell$, then $\left(
\lambda_{1},\lambda_{2},\ldots,\lambda_{\ell}\right)  $ is a
partition\footnote{because of Observation 1 \textbf{(d)}} with largest part
$\leq n-k$\ \ \ \ \footnote{since all parts $\lambda_{1},\lambda_{2}%
,\ldots,\lambda_{\ell}$ of this partition are $\leq n-k$ (because
$n-k\geq\lambda_{1}\geq\lambda_{2}\geq\cdots\geq\lambda_{k}$)} and length
$\ell$. Thus, we obtain a map%
\begin{align*}
&  \text{from }\left\{  k\text{-tuples }\left(  \lambda_{1},\lambda_{2}%
,\ldots,\lambda_{k}\right)  \in\mathbb{N}^{k}\text{ satisfying }n-k\geq
\lambda_{1}\geq\lambda_{2}\geq\cdots\geq\lambda_{k}\geq0\right. \\
&  \ \ \ \ \ \ \ \ \ \ \ \ \ \ \ \ \ \ \ \ \left.  \text{ and }%
\operatorname*{numpos}\left(  \lambda_{1},\lambda_{2},\ldots,\lambda
_{k}\right)  =\ell\right\} \\
&  \text{to }\left\{  \text{partitions with largest part }\leq n-k\text{ and
length }\ell\right\}
\end{align*}
which sends any $k$-tuple $\left(  \lambda_{1},\lambda_{2},\ldots,\lambda
_{k}\right)  $ to the partition $\left(  \lambda_{1},\lambda_{2}%
,\ldots,\lambda_{\ell}\right)  $. On the other hand, we have a map%
\begin{align*}
&  \text{from }\left\{  \text{partitions with largest part }\leq n-k\text{ and
length }\ell\right\} \\
&  \text{to }\left\{  k\text{-tuples }\left(  \lambda_{1},\lambda_{2}%
,\ldots,\lambda_{k}\right)  \in\mathbb{N}^{k}\text{ satisfying }n-k\geq
\lambda_{1}\geq\lambda_{2}\geq\cdots\geq\lambda_{k}\geq0\right. \\
&  \ \ \ \ \ \ \ \ \ \ \ \ \ \ \ \ \ \ \ \ \left.  \text{ and }%
\operatorname*{numpos}\left(  \lambda_{1},\lambda_{2},\ldots,\lambda
_{k}\right)  =\ell\right\}
\end{align*}
which sends any partition $\left(  \lambda_{1},\lambda_{2},\ldots
,\lambda_{\ell}\right)  $ to the $k$-tuple $\left(  \lambda_{1},\lambda
_{2},\ldots,\lambda_{\ell},\underbrace{0,0,\ldots,0}_{k-\ell\text{ zeroes}%
}\right)  $\ \ \ \ \footnote{because we have shown in the previous paragraph
that if $\left(  \lambda_{1},\lambda_{2},\ldots,\lambda_{\ell}\right)  $ is a
partition with largest part $\leq n-k$, then we can extend it to a $k$-tuple
$\left(  \lambda_{1},\lambda_{2},\ldots,\lambda_{k}\right)  =\left(
\lambda_{1},\lambda_{2},\ldots,\lambda_{\ell},\underbrace{0,0,\ldots
,0}_{k-\ell\text{ zeroes}}\right)  $ by inserting $k-\ell$ zeroes at the end,
and this extended $k$-tuple will satisfy $n-k\geq\lambda_{1}\geq\lambda
_{2}\geq\cdots\geq\lambda_{k}\geq0$ and $\operatorname*{numpos}\left(
\lambda_{1},\lambda_{2},\ldots,\lambda_{k}\right)  =\ell$}. These two maps are
mutually inverse\footnote{Indeed, the first map removes the last $k-\ell$
entries from a $k$-tuple, whereas the second map inserts $k-\ell$ zeroes at
the end of the partition. Thus, if we apply the first map after the second
map, we clearly recover the partition that we started with. If we apply the
second map after the first map, then we end up replacing the last $k-\ell$
entries of our $k$-tuple by zeroes. However, if $\left(  \lambda_{1}%
,\lambda_{2},\ldots,\lambda_{k}\right)  \in\mathbb{N}^{k}$ is any $k$-tuple
satisfying $n-k\geq\lambda_{1}\geq\lambda_{2}\geq\cdots\geq\lambda_{k}\geq0$
and $\operatorname*{numpos}\left(  \lambda_{1},\lambda_{2},\ldots,\lambda
_{k}\right)  =\ell$, then the last $k-\ell$ entries of this $k$-tuple are $0$
(by Observation 1 \textbf{(b)}), and therefore the $k$-tuple does not change
if we replace these $k-\ell$ entries by zeroes.}, and therefore are
bijections. In particular, this shows that the first map is a bijection. This
bijection allows us to replace our partitions $\left(  \lambda_{1},\lambda
_{2},\ldots,\lambda_{\ell}\right)  \in\mathbb{N}^{\ell}$ by $k$-tuples
$\left(  \lambda_{1},\lambda_{2},\ldots,\lambda_{k}\right)  \in\mathbb{N}^{k}$
in the sum $\sum_{\substack{\left(  \lambda_{1},\lambda_{2},\ldots
,\lambda_{\ell}\right)  \in\mathbb{N}^{\ell};\\n-k\geq\lambda_{1}\geq
\lambda_{2}\geq\cdots\geq\lambda_{\ell}>0}}q^{\lambda_{1}+\lambda_{2}%
+\cdots+\lambda_{\ell}}$; we thus find%
\begin{align*}
\sum_{\substack{\left(  \lambda_{1},\lambda_{2},\ldots,\lambda_{\ell}\right)
\in\mathbb{N}^{\ell};\\n-k\geq\lambda_{1}\geq\lambda_{2}\geq\cdots\geq
\lambda_{\ell}>0}}q^{\lambda_{1}+\lambda_{2}+\cdots+\lambda_{\ell}}  &
=\sum_{\substack{\left(  \lambda_{1},\lambda_{2},\ldots,\lambda_{k}\right)
\in\mathbb{N}^{k};\\n-k\geq\lambda_{1}\geq\lambda_{2}\geq\cdots\geq\lambda
_{k}\geq0;\\\operatorname*{numpos}\left(  \lambda_{1},\lambda_{2}%
,\ldots,\lambda_{k}\right)  =\ell}}\ \ \underbrace{q^{\lambda_{1}+\lambda
_{2}+\cdots+\lambda_{\ell}}}_{\substack{=q^{\lambda_{1}+\lambda_{2}%
+\cdots+\lambda_{k}}\\\text{(by Observation 1 \textbf{(c)})}}}\\
&  =\sum_{\substack{\left(  \lambda_{1},\lambda_{2},\ldots,\lambda_{k}\right)
\in\mathbb{N}^{k};\\n-k\geq\lambda_{1}\geq\lambda_{2}\geq\cdots\geq\lambda
_{k}\geq0;\\\operatorname*{numpos}\left(  \lambda_{1},\lambda_{2}%
,\ldots,\lambda_{k}\right)  =\ell}}q^{\lambda_{1}+\lambda_{2}+\cdots
+\lambda_{k}}.
\end{align*}


Now, (\ref{pf.prop.pars.qbinom.alt-defs.a.2}) becomes%
\begin{align}
\sum_{\substack{\lambda\text{ is a partition}\\\text{with largest part }\leq
n-k\\\text{and length }\ell}}q^{\left\vert \lambda\right\vert }  &
=\sum_{\substack{\left(  \lambda_{1},\lambda_{2},\ldots,\lambda_{\ell}\right)
\in\mathbb{N}^{\ell};\\n-k\geq\lambda_{1}\geq\lambda_{2}\geq\cdots\geq
\lambda_{\ell}>0}}q^{\lambda_{1}+\lambda_{2}+\cdots+\lambda_{\ell}}\nonumber\\
&  =\sum_{\substack{\left(  \lambda_{1},\lambda_{2},\ldots,\lambda_{k}\right)
\in\mathbb{N}^{k};\\n-k\geq\lambda_{1}\geq\lambda_{2}\geq\cdots\geq\lambda
_{k}\geq0;\\\operatorname*{numpos}\left(  \lambda_{1},\lambda_{2}%
,\ldots,\lambda_{k}\right)  =\ell}}q^{\lambda_{1}+\lambda_{2}+\cdots
+\lambda_{k}}. \label{pf.prop.pars.qbinom.alt-defs.a.5}%
\end{align}


Now, forget that we fixed $\ell$. We thus have proved
(\ref{pf.prop.pars.qbinom.alt-defs.a.5}) for each $\ell\in\left\{
0,1,\ldots,k\right\}  $. Now, (\ref{pf.prop.pars.qbinom.alt-defs.a.1}) becomes%
\begin{align*}
\dbinom{n}{k}_{q}  &  =\sum_{\ell=0}^{k}\ \ \underbrace{\sum
_{\substack{\lambda\text{ is a partition}\\\text{with largest part }\leq
n-k\\\text{and length }\ell}}q^{\left\vert \lambda\right\vert }}%
_{\substack{=\sum_{\substack{\left(  \lambda_{1},\lambda_{2},\ldots
,\lambda_{k}\right)  \in\mathbb{N}^{k};\\n-k\geq\lambda_{1}\geq\lambda_{2}%
\geq\cdots\geq\lambda_{k}\geq0;\\\operatorname*{numpos}\left(  \lambda
_{1},\lambda_{2},\ldots,\lambda_{k}\right)  =\ell}}q^{\lambda_{1}+\lambda
_{2}+\cdots+\lambda_{k}}\\\text{(by (\ref{pf.prop.pars.qbinom.alt-defs.a.5}%
))}}}\\
&  =\underbrace{\sum_{\ell=0}^{k}\ \ \sum_{\substack{\left(  \lambda
_{1},\lambda_{2},\ldots,\lambda_{k}\right)  \in\mathbb{N}^{k};\\n-k\geq
\lambda_{1}\geq\lambda_{2}\geq\cdots\geq\lambda_{k}\geq
0;\\\operatorname*{numpos}\left(  \lambda_{1},\lambda_{2},\ldots,\lambda
_{k}\right)  =\ell}}}_{=\sum_{\substack{\left(  \lambda_{1},\lambda_{2}%
,\ldots,\lambda_{k}\right)  \in\mathbb{N}^{k};\\n-k\geq\lambda_{1}\geq
\lambda_{2}\geq\cdots\geq\lambda_{k}\geq0}}}q^{\lambda_{1}+\lambda_{2}%
+\cdots+\lambda_{k}}\\
&  =\sum_{\substack{\left(  \lambda_{1},\lambda_{2},\ldots,\lambda_{k}\right)
\in\mathbb{N}^{k};\\n-k\geq\lambda_{1}\geq\lambda_{2}\geq\cdots\geq\lambda
_{k}\geq0}}q^{\lambda_{1}+\lambda_{2}+\cdots+\lambda_{k}}=\sum
_{\substack{\left(  i_{1},i_{2},\ldots,i_{k}\right)  \in\mathbb{N}^{k};\\0\leq
i_{1}\leq i_{2}\leq\cdots\leq i_{k}\leq n-k}}\underbrace{q^{i_{k}%
+i_{k-1}+\cdots+i_{1}}}_{=q^{i_{1}+i_{2}+\cdots+i_{k}}}\\
&  \ \ \ \ \ \ \ \ \ \ \left(
\begin{array}
[c]{c}%
\text{here, we have reversed the }k\text{-tuple }\left(  \lambda_{1}%
,\lambda_{2},\ldots,\lambda_{k}\right)  \text{,}\\
\text{i.e., we have substituted }\left(  i_{k},i_{k-1},\ldots,i_{1}\right)
\text{ for }\left(  \lambda_{1},\lambda_{2},\ldots,\lambda_{k}\right) \\
\text{in our sum}%
\end{array}
\right) \\
&  =\sum_{\substack{\left(  i_{1},i_{2},\ldots,i_{k}\right)  \in\mathbb{N}%
^{k};\\0\leq i_{1}\leq i_{2}\leq\cdots\leq i_{k}\leq n-k}}q^{i_{1}%
+i_{2}+\cdots+i_{k}}=\sum_{0\leq i_{1}\leq i_{2}\leq\cdots\leq i_{k}\leq
n-k}q^{i_{1}+i_{2}+\cdots+i_{k}}.
\end{align*}
This proves Proposition \ref{prop.pars.qbinom.alt-defs} \textbf{(a)}. \medskip

\textbf{(b)} There is a bijection%
\begin{align*}
&  \text{from }\left\{  \text{weakly increasing }k\text{-tuples }\left(
i_{1},i_{2},\ldots,i_{k}\right)  \in\left\{  0,1,\ldots,n-k\right\}
^{k}\right\} \\
&  \text{to }\left\{  \text{strictly increasing }k\text{-tuples }\left(
s_{1},s_{2},\ldots,s_{k}\right)  \in\left\{  1,2,\ldots,n\right\}
^{k}\right\}
\end{align*}
that sends each weakly increasing $k$-tuple $\left(  i_{1},i_{2},\ldots
,i_{k}\right)  $ to $\left(  i_{1}+1,\ i_{2}+2,\ \ldots,\ i_{k}+k\right)  $
(you can think of it as \textquotedblleft spacing the $i_{j}$s
apart\textquotedblright, i.e., increasing the distance between any two
consecutive $i_{j}$'s by $1$ and also increasing $i_{1}$ by $1$). The inverse
of this bijection sends each strictly increasing $k$-tuple $\left(
s_{1},s_{2},\ldots,s_{k}\right)  $ to $\left(  s_{1}-1,\ s_{2}-2,\ \ldots
,\ s_{k}-k\right)  $. Thus, we can substitute $\left(  s_{1}-1,\ s_{2}%
-2,\ \ldots,\ s_{k}-k\right)  $ for $\left(  i_{1},i_{2},\ldots,i_{k}\right)
$ in the sum%
\[
\sum_{0\leq i_{1}\leq i_{2}\leq\cdots\leq i_{k}\leq n-k}q^{i_{1}+i_{2}%
+\cdots+i_{k}}.
\]
Hence we obtain%
\begin{align*}
\sum_{0\leq i_{1}\leq i_{2}\leq\cdots\leq i_{k}\leq n-k}q^{i_{1}+i_{2}%
+\cdots+i_{k}}  &  =\sum_{1\leq s_{1}<s_{2}<\cdots<s_{k}\leq n}%
\underbrace{q^{\left(  s_{1}-1\right)  +\left(  s_{2}-2\right)  +\cdots
+\left(  s_{k}-k\right)  }}_{=q^{\left(  s_{1}+s_{2}+\cdots+s_{k}\right)
-\left(  1+2+\cdots+k\right)  }}\\
&  =\sum_{1\leq s_{1}<s_{2}<\cdots<s_{k}\leq n}q^{\left(  s_{1}+s_{2}%
+\cdots+s_{k}\right)  -\left(  1+2+\cdots+k\right)  }.
\end{align*}


On the other hand, there is a bijection%
\begin{align*}
&  \text{from }\left\{  \text{strictly increasing }k\text{-tuples }\left(
s_{1},s_{2},\ldots,s_{k}\right)  \in\left\{  1,2,\ldots,n\right\}
^{k}\right\} \\
&  \text{to }\left\{  k\text{-element subsets of }\left\{  1,2,\ldots
,n\right\}  \right\}
\end{align*}
that sends each $k$-tuple $\left(  s_{1},s_{2},\ldots,s_{k}\right)  $ to the
subset $\left\{  s_{1},s_{2},\ldots,s_{k}\right\}  $. (This map is indeed a
bijection, because any $k$-element subset of $\left\{  1,2,\ldots,n\right\}  $
can be written as $\left\{  s_{1},s_{2},\ldots,s_{k}\right\}  $ for a unique
strictly increasing $k$-tuple $\left(  s_{1},s_{2},\ldots,s_{k}\right)
\in\left\{  1,2,\ldots,n\right\}  ^{k}$; in fact, this is simply saying that
there is a unique way of listing the elements of this subset in increasing order.)

Because of this bijection, we have%
\[
\sum_{1\leq s_{1}<s_{2}<\cdots<s_{k}\leq n}q^{\left(  s_{1}+s_{2}+\cdots
+s_{k}\right)  -\left(  1+2+\cdots+k\right)  }=\sum_{\substack{S\subseteq
\left\{  1,2,\ldots,n\right\}  ;\\\left\vert S\right\vert =k}%
}q^{\operatorname*{sum}S-\left(  1+2+\cdots+k\right)  }%
\]
(because $s_{1}+s_{2}+\cdots+s_{k}=\operatorname*{sum}\left\{  s_{1}%
,s_{2},\ldots,s_{k}\right\}  $ for any strictly increasing $k$-tuple $\left(
s_{1},s_{2},\ldots,s_{k}\right)  \in\left\{  1,2,\ldots,n\right\}  ^{k}$).

Now, Proposition \ref{prop.pars.qbinom.alt-defs} \textbf{(a)} yields%
\begin{align*}
\dbinom{n}{k}_{q}  &  =\sum_{0\leq i_{1}\leq i_{2}\leq\cdots\leq i_{k}\leq
n-k}q^{i_{1}+i_{2}+\cdots+i_{k}}=\sum_{1\leq s_{1}<s_{2}<\cdots<s_{k}\leq
n}q^{\left(  s_{1}+s_{2}+\cdots+s_{k}\right)  -\left(  1+2+\cdots+k\right)
}\\
&  =\sum_{\substack{S\subseteq\left\{  1,2,\ldots,n\right\}  ;\\\left\vert
S\right\vert =k}}q^{\operatorname*{sum}S-\left(  1+2+\cdots+k\right)  }.
\end{align*}
This proves Proposition \ref{prop.pars.qbinom.alt-defs} \textbf{(b)}. \medskip

\textbf{(c)} Proposition \ref{prop.pars.qbinom.alt-defs} \textbf{(b)} yields
\[
\dbinom{n}{k}_{q}=\sum_{\substack{S\subseteq\left\{  1,2,\ldots,n\right\}
;\\\left\vert S\right\vert =k}}q^{\operatorname*{sum}S-\left(  1+2+\cdots
+k\right)  }.
\]
Substituting $1$ for $q$ in this equality, we find%
\begin{align*}
\dbinom{n}{k}_{1}  &  =\sum_{\substack{S\subseteq\left\{  1,2,\ldots
,n\right\}  ;\\\left\vert S\right\vert =k}}\underbrace{1^{\operatorname*{sum}%
S-\left(  1+2+\cdots+k\right)  }}_{=1}=\sum_{\substack{S\subseteq\left\{
1,2,\ldots,n\right\}  ;\\\left\vert S\right\vert =k}}1\\
&  =\left(  \text{\# of subsets }S\subseteq\left\{  1,2,\ldots,n\right\}
\text{ satisfying }\left\vert S\right\vert =k\right) \\
&  =\left(  \text{\# of }k\text{-element subsets of }\left\{  1,2,\ldots
,n\right\}  \right)  =\dbinom{n}{k}.
\end{align*}
This proves Proposition \ref{prop.pars.qbinom.alt-defs} \textbf{(c)}.
\end{proof}

The following property of $q$-binomial coefficients generalizes Proposition
\ref{prop.binom.0}:

\begin{proposition}
\label{prop.pars.qbinom.0}Let $n,k\in\mathbb{N}$ satisfy $k>n$. Then,
$\dbinom{n}{k}_{q}=0$.
\end{proposition}

\begin{proof}
From $k>n$, we obtain $n-k<0$. The definition of $\dbinom{n}{k}_{q}$ yields%
\begin{equation}
\dbinom{n}{k}_{q}=\sum_{\substack{\lambda\text{ is a partition}\\\text{with
largest part }\leq n-k\\\text{and length }\leq k}}q^{\left\vert \lambda
\right\vert }. \label{pf.prop.pars.qbinom.0.1}%
\end{equation}
The sum on the right hand side is an empty sum, since there exists no
partition with largest part $\leq n-k$ (because $n-k<0$). Thus,
(\ref{pf.prop.pars.qbinom.0.1}) rewrites as $\dbinom{n}{k}_{q}=\left(
\text{empty sum}\right)  =0$, and this proves Proposition
\ref{prop.pars.qbinom.0}.
\end{proof}

\begin{proposition}
\label{prop.pars.qbinom.n0}We have $\dbinom{n}{0}_{q}=\dbinom{n}{n}_{q}=1$ for
each $n\in\mathbb{N}$.
\end{proposition}

\begin{proof}
This is easy and left as a homework exercise (Exercise
\ref{exe.pars.qbinom.basics} \textbf{(a)}).
\end{proof}

The next convention mirrors a convention we made for the (usual) binomial coefficients:

\begin{convention}
\label{conv.pars.qbinom.neg-k}Let $n\in\mathbb{N}$. For any $k\notin%
\mathbb{Z}$, we set $\dbinom{n}{k}_{q}:=0$.
\end{convention}

The following theorem gives not one, but two analogues (\textquotedblleft%
$q$-analogues\textquotedblright) of the recurrence relation $\dbinom{n}%
{k}=\dbinom{n-1}{k-1}+\dbinom{n-1}{k}$ (from Proposition \ref{prop.binom.rec}):

\begin{theorem}
\label{thm.pars.qbinom.rec}Let $n$ be a positive integer. Let $k\in\mathbb{N}%
$. Then: \medskip

\textbf{(a)} We have%
\[
\dbinom{n}{k}_{q}=q^{n-k}\dbinom{n-1}{k-1}_{q}+\dbinom{n-1}{k}_{q}.
\]


\textbf{(b)} We have%
\[
\dbinom{n}{k}_{q}=\dbinom{n-1}{k-1}_{q}+q^{k}\dbinom{n-1}{k}_{q}.
\]

\end{theorem}

\begin{proof}
\textbf{(a)} This is similar to the combinatorial proof of the recurrence
relation for binomial coefficients.

If $k=0$, then the claim we are proving boils down to $1=q^{n-k}0+1$ (because
Proposition \ref{prop.pars.qbinom.n0} yields $\dbinom{n}{0}_{q}=1$ and
$\dbinom{n-1}{0}_{q}=1$, and because Convention \ref{conv.pars.qbinom.neg-k}
yields $\dbinom{n-1}{-1}_{q}=0$). Hence, we WLOG assume that $k>0$. Thus,
$k-1\in\mathbb{N}$.

Proposition \ref{prop.pars.qbinom.alt-defs} \textbf{(b)} says that%
\begin{equation}
\dbinom{n}{k}_{q}=\sum_{\substack{S\subseteq\left\{  1,2,\ldots,n\right\}
;\\\left\vert S\right\vert =k}}q^{\operatorname*{sum}S-\left(  1+2+\cdots
+k\right)  }. \label{pf.thm.pars.qbinom.rec.nk}%
\end{equation}
Proposition \ref{prop.pars.qbinom.alt-defs} \textbf{(b)} (applied to $n-1$ and
$k-1$ instead of $n$ and $k$) yields%
\begin{equation}
\dbinom{n-1}{k-1}_{q}=\sum_{\substack{S\subseteq\left\{  1,2,\ldots
,n-1\right\}  ;\\\left\vert S\right\vert =k-1}}q^{\operatorname*{sum}S-\left(
1+2+\cdots+\left(  k-1\right)  \right)  }.
\label{pf.thm.pars.qbinom.rec.n-1k-1}%
\end{equation}
Proposition \ref{prop.pars.qbinom.alt-defs} \textbf{(b)} (applied to $n-1$
instead of $n$) yields%
\begin{equation}
\dbinom{n-1}{k}_{q}=\sum_{\substack{S\subseteq\left\{  1,2,\ldots,n-1\right\}
;\\\left\vert S\right\vert =k}}q^{\operatorname*{sum}S-\left(  1+2+\cdots
+k\right)  }. \label{pf.thm.pars.qbinom.rec.n-1k}%
\end{equation}


Let us now make two definitions:

\begin{itemize}
\item A \emph{type-1 subset} will mean a $k$-element subset of $\left\{
1,2,\ldots,n\right\}  $ that contains $n$;

\item A \emph{type-2 subset} will mean a $k$-element subset of $\left\{
1,2,\ldots,n\right\}  $ that does not contain $n$.
\end{itemize}

Each $k$-element subset of $\left\{  1,2,\ldots,n\right\}  $ is either type-1
or type-2 (but not both at the same time). Thus,%
\begin{align*}
&  \sum_{\substack{S\subseteq\left\{  1,2,\ldots,n\right\}  ;\\\left\vert
S\right\vert =k}}q^{\operatorname*{sum}S-\left(  1+2+\cdots+k\right)  }\\
&  =\sum_{\substack{S\subseteq\left\{  1,2,\ldots,n\right\}  ;\\\left\vert
S\right\vert =k;\\S\text{ is type-1}}}q^{\operatorname*{sum}S-\left(
1+2+\cdots+k\right)  }+\sum_{\substack{S\subseteq\left\{  1,2,\ldots
,n\right\}  ;\\\left\vert S\right\vert =k;\\S\text{ is type-2}}%
}q^{\operatorname*{sum}S-\left(  1+2+\cdots+k\right)  }.
\end{align*}


The type-2 subsets are precisely the $k$-element subsets of $\left\{
1,2,\ldots,n-1\right\}  $. Hence,%
\[
\sum_{\substack{S\subseteq\left\{  1,2,\ldots,n\right\}  ;\\\left\vert
S\right\vert =k;\\S\text{ is type-2}}}q^{\operatorname*{sum}S-\left(
1+2+\cdots+k\right)  }=\sum_{\substack{S\subseteq\left\{  1,2,\ldots
,n-1\right\}  ;\\\left\vert S\right\vert =k}}q^{\operatorname*{sum}S-\left(
1+2+\cdots+k\right)  }=\dbinom{n-1}{k}_{q}%
\]
(by (\ref{pf.thm.pars.qbinom.rec.n-1k})).

The type-1 subsets are just the $\left(  k-1\right)  $-element subsets of
$\left\{  1,2,\ldots,n-1\right\}  $ with an $n$ inserted into them; i.e., the
map%
\begin{align*}
\left\{  \left(  k-1\right)  \text{-element subsets of }\left\{
1,2,\ldots,n-1\right\}  \right\}   &  \rightarrow\left\{  \text{type-1
subsets}\right\}  ,\\
S  &  \mapsto S\cup\left\{  n\right\}
\end{align*}
is a bijection. Hence, substituting $S\cup\left\{  n\right\}  $ for $S$ in the
sum, we find%
\begin{align*}
\sum_{\substack{S\subseteq\left\{  1,2,\ldots,n\right\}  ;\\\left\vert
S\right\vert =k;\\S\text{ is type-1}}}q^{\operatorname*{sum}S-\left(
1+2+\cdots+k\right)  }  &  =\sum_{\substack{S\subseteq\left\{  1,2,\ldots
,n-1\right\}  ;\\\left\vert S\right\vert =k-1}%
}\underbrace{q^{\operatorname*{sum}\left(  S\cup\left\{  n\right\}  \right)
-\left(  1+2+\cdots+k\right)  }}_{\substack{=q^{\operatorname*{sum}S+n-\left(
1+2+\cdots+k\right)  }\\\text{(since }S\subseteq\left\{  1,2,\ldots
,n-1\right\}  \text{ entails }n\notin S\\\text{and thus }\operatorname*{sum}%
\left(  S\cup\left\{  n\right\}  \right)  =\operatorname*{sum}S+n\text{)}}}\\
&  =\sum_{\substack{S\subseteq\left\{  1,2,\ldots,n-1\right\}  ;\\\left\vert
S\right\vert =k-1}}\underbrace{q^{\operatorname*{sum}S+n-\left(
1+2+\cdots+k\right)  }}_{\substack{=q^{\operatorname*{sum}S+n-\left(
1+2+\cdots+\left(  k-1\right)  \right)  -k}\\=q^{n-k}q^{\operatorname*{sum}%
S-\left(  1+2+\cdots+\left(  k-1\right)  \right)  }}}\\
&  =\sum_{\substack{S\subseteq\left\{  1,2,\ldots,n-1\right\}  ;\\\left\vert
S\right\vert =k-1}}q^{n-k}q^{\operatorname*{sum}S-\left(  1+2+\cdots+\left(
k-1\right)  \right)  }\\
&  =q^{n-k}\underbrace{\sum_{\substack{S\subseteq\left\{  1,2,\ldots
,n-1\right\}  ;\\\left\vert S\right\vert =k-1}}q^{\operatorname*{sum}S-\left(
1+2+\cdots+\left(  k-1\right)  \right)  }}_{\substack{=\dbinom{n-1}{k-1}%
_{q}\\\text{(by (\ref{pf.thm.pars.qbinom.rec.n-1k-1}))}}}\\
&  =q^{n-k}\dbinom{n-1}{k-1}_{q}.
\end{align*}


All that's left to do is combining what we have found:%
\begin{align*}
\dbinom{n}{k}_{q}  &  =\sum_{\substack{S\subseteq\left\{  1,2,\ldots
,n\right\}  ;\\\left\vert S\right\vert =k}}q^{\operatorname*{sum}S-\left(
1+2+\cdots+k\right)  }\ \ \ \ \ \ \ \ \ \ \left(  \text{by
(\ref{pf.thm.pars.qbinom.rec.nk})}\right) \\
&  =\underbrace{\sum_{\substack{S\subseteq\left\{  1,2,\ldots,n\right\}
;\\\left\vert S\right\vert =k;\\S\text{ is type-1}}}q^{\operatorname*{sum}%
S-\left(  1+2+\cdots+k\right)  }}_{=q^{n-k}\dbinom{n-1}{k-1}_{q}%
}+\underbrace{\sum_{\substack{S\subseteq\left\{  1,2,\ldots,n\right\}
;\\\left\vert S\right\vert =k;\\S\text{ is type-2}}}q^{\operatorname*{sum}%
S-\left(  1+2+\cdots+k\right)  }}_{=\dbinom{n-1}{k}_{q}}\\
&  =q^{n-k}\dbinom{n-1}{k-1}_{q}+\dbinom{n-1}{k}_{q}.
\end{align*}
This proves Theorem \ref{thm.pars.qbinom.rec} \textbf{(a)}. \medskip

\textbf{(b)} This is somewhat similar to Theorem \ref{thm.pars.qbinom.rec}
\textbf{(a)} (but a little bit more complicated). It is left as a homework
exercise (Exercise \ref{exe.pars.qbinom.basics} \textbf{(b)}).
\end{proof}

Next, we shall derive a $q$-analogue of the formula $\dbinom{n}{k}%
=\dfrac{n\left(  n-1\right)  \cdots\left(  n-k+1\right)  }{k!}=\dfrac{n\left(
n-1\right)  \cdots\left(  n-k+1\right)  }{k\left(  k-1\right)  \cdots1}$:

\begin{theorem}
\label{thm.pars.qbinom.quot1}Let $n,k\in\mathbb{N}$ satisfy $n\geq k$. Then:
\medskip

\textbf{(a)} We have%
\begin{align*}
&  \left(  1-q^{k}\right)  \left(  1-q^{k-1}\right)  \cdots\left(
1-q^{1}\right)  \cdot\dbinom{n}{k}_{q}\\
&  =\left(  1-q^{n}\right)  \left(  1-q^{n-1}\right)  \cdots\left(
1-q^{n-k+1}\right)  .
\end{align*}


\textbf{(b)} We have%
\[
\dbinom{n}{k}_{q}=\dfrac{\left(  1-q^{n}\right)  \left(  1-q^{n-1}\right)
\cdots\left(  1-q^{n-k+1}\right)  }{\left(  1-q^{k}\right)  \left(
1-q^{k-1}\right)  \cdots\left(  1-q^{1}\right)  }%
\]
(in the ring $\mathbb{Z}\left[  \left[  q\right]  \right]  $ or in the field
of rational functions over $\mathbb{Q}$).
\end{theorem}

Note that part \textbf{(b)} of Theorem \ref{thm.pars.qbinom.quot1} is the more
intuitive statement, but part \textbf{(a)} is easier to substitute things in
(because substituting something for $q$ in part \textbf{(b)} requires showing
that the denominator remains invertible, whereas part \textbf{(a)} has no
denominators and thus requires no such diligence).

\begin{proof}
[Proof of Theorem \ref{thm.pars.qbinom.quot1}.]This is left as a homework
exercise (Exercise \ref{exe.pars.qbinom.basics} \textbf{(c)}). (Use induction
on $n$ and Theorem \ref{thm.pars.qbinom.rec}.)
\end{proof}

\begin{remark}
If I just gave you the fraction $\dfrac{\left(  1-q^{n}\right)  \left(
1-q^{n-1}\right)  \cdots\left(  1-q^{n-k+1}\right)  }{\left(  1-q^{k}\right)
\left(  1-q^{k-1}\right)  \cdots\left(  1-q^{1}\right)  }$, you would be
surprised to hear that it is a polynomial (i.e., that the denominator divides
the numerator) and has nonnegative coefficients. But given the way we defined
$\dbinom{n}{k}_{q}$, you are now getting this for free from Theorem
\ref{thm.pars.qbinom.quot1}.
\end{remark}

Theorem \ref{thm.pars.qbinom.quot1} \textbf{(b)} can be rewritten in a
somewhat simpler way using the following notations:

\begin{definition}
\label{def.pars.qbinom.qint}\textbf{(a)} For any $n\in\mathbb{N}$, define the
$q$\emph{-integer }$\left[  n\right]  _{q}$ to be
\[
\left[  n\right]  _{q}:=q^{0}+q^{1}+\cdots+q^{n-1}\in\mathbb{Z}\left[
q\right]  .
\]


\textbf{(b)} For any $n\in\mathbb{N}$, define the $q$\emph{-factorial
}$\left[  n\right]  _{q}!$ to be%
\[
\left[  n\right]  _{q}!\ :=\left[  1\right]  _{q}\left[  2\right]  _{q}%
\cdots\left[  n\right]  _{q}\in\mathbb{Z}\left[  q\right]  .
\]


\textbf{(c)} As usual, if $a$ is an element of a ring $A$, then $\left[
n\right]  _{a}$ and $\left[  n\right]  _{a}!$ will mean the results of
substituting $a$ for $q$ in $\left[  n\right]  _{q}$ and $\left[  n\right]
_{q}!$, respectively. Thus, explicitly, $\left[  n\right]  _{a}=a^{0}%
+a^{1}+\cdots+a^{n-1}$ and $\left[  n\right]  _{a}!=\left[  1\right]
_{a}\left[  2\right]  _{a}\cdots\left[  n\right]  _{a}$.
\end{definition}

Alternative notations for $\left[  n\right]  _{q}!$ are $\left[  n\right]
!_{q}$ (used in \cite{Loehr-BC}) and the boldface $\mathbf{(}\boldsymbol{n}%
\mathbf{)!}$ (used in \cite{Stanley-EC1}).

\begin{remark}
\label{rmk.pars.qbinom.qint.frac}For any $n\in\mathbb{N}$, we have%
\[
\left[  n\right]  _{q}=\dfrac{1-q^{n}}{1-q}\ \ \ \ \ \ \ \ \ \ \left(
\text{in }\mathbb{Z}\left[  \left[  q\right]  \right]  \text{ or in the ring
of rational functions over }\mathbb{Q}\right)
\]
and%
\[
\left[  n\right]  _{1}=n\ \ \ \ \ \ \ \ \ \ \text{and}%
\ \ \ \ \ \ \ \ \ \ \left[  n\right]  _{1}!=n!.
\]

\end{remark}

\begin{proof}
[Proof of Remark \ref{rmk.pars.qbinom.qint.frac}.]Let $n\in\mathbb{N}$. We
have%
\[
\left[  n\right]  _{q}:=q^{0}+q^{1}+\cdots+q^{n-1}=\dfrac{1-q^{n}}{1-q},
\]
since%
\begin{align*}
\left(  1-q\right)  \left(  q^{0}+q^{1}+\cdots+q^{n-1}\right)   &  =\left(
q^{0}+q^{1}+\cdots+q^{n-1}\right)  -q\left(  q^{0}+q^{1}+\cdots+q^{n-1}\right)
\\
&  =\left(  q^{0}+q^{1}+\cdots+q^{n-1}\right)  -\left(  q^{1}+q^{2}%
+\cdots+q^{n}\right) \\
&  =\underbrace{q^{0}}_{=1}-\,q^{n}=1-q^{n}.
\end{align*}


Furthermore, substituting $1$ for $q$ in the equality $\left[  n\right]
_{q}=q^{0}+q^{1}+\cdots+q^{n-1}$, we obtain%
\begin{equation}
\left[  n\right]  _{1}=1^{0}+1^{1}+\cdots+1^{n-1}=\underbrace{1+1+\cdots
+1}_{n\text{ times}}=n. \label{pf.rmk.pars.qbinom.qint.frac.4}%
\end{equation}
Substituting $1$ for $q$ in the equality $\left[  n\right]  _{q}!=\left[
1\right]  _{q}\left[  2\right]  _{q}\cdots\left[  n\right]  _{q}$, we obtain%
\[
\left[  n\right]  _{1}!=\underbrace{\left[  1\right]  _{1}}%
_{\substack{=1\\\text{(by (\ref{pf.rmk.pars.qbinom.qint.frac.4}))}%
}}\underbrace{\left[  2\right]  _{1}}_{\substack{=2\\\text{(by
(\ref{pf.rmk.pars.qbinom.qint.frac.4}))}}}\cdots\underbrace{\left[  n\right]
_{1}}_{\substack{=n\\\text{(by (\ref{pf.rmk.pars.qbinom.qint.frac.4}))}%
}}=1\cdot2\cdot\cdots\cdot n=n!.
\]

\end{proof}

\begin{theorem}
\label{thm.pars.qbinom.quot2}Let $n,k\in\mathbb{N}$ with $n\geq k$. Then,%
\[
\dbinom{n}{k}_{q}=\dfrac{\left[  n\right]  _{q}\left[  n-1\right]  _{q}%
\cdots\left[  n-k+1\right]  _{q}}{\left[  k\right]  _{q}!}=\dfrac{\left[
n\right]  _{q}!}{\left[  k\right]  _{q}!\cdot\left[  n-k\right]  _{q}!}%
\]
(in the ring $\mathbb{Z}\left[  \left[  q\right]  \right]  $ or in the ring of
rational functions over $\mathbb{Q}$).
\end{theorem}

\begin{proof}
This is left as a homework exercise (Exercise \ref{exe.pars.qbinom.basics}
\textbf{(d)}).
\end{proof}

A consequence of this theorem is the following symmetry property of
$q$-binomial coefficients:

\begin{proposition}
\label{prop.pars.qbinom.symm}Let $n,k\in\mathbb{N}$. Then,%
\[
\dbinom{n}{k}_{q}=\dbinom{n}{n-k}_{q}.
\]

\end{proposition}

\begin{proof}
This is left as a homework exercise (Exercise \ref{exe.pars.qbinom.basics}
\textbf{(e)}).
\end{proof}

